\section{Experimental Setup}
\label{sec:setup}

\subsection{Pre-registration}

Every experiment in this paper was pre-registered: the acceptance gates,
the selected operating points, the statistical procedure and the seed
policy were committed to writing \emph{before} the experiment ran, and
each amendment carries a timestamp and a stated reason. Amendments made
after a result was visible are not permitted, and none were made.

This matters for how the paper should be read. A gate that failed is
reported as failed rather than reinterpreted, and no parameter was
adjusted in response to any result reported here. The design evolution in
Section~\ref{sec:results} contains two variants that failed their gates;
neither differs from the working variant by any parameter value.

\subsection{Methods compared}

\begin{itemize}
\item \textbf{P10} --- periodic, $10$~Hz per channel.
\item \textbf{P20} --- periodic, $20$~Hz per channel.
\item \textbf{State-event} --- conventional state-error
      trigger~\cite{tabuada2007event,dimarogonas2012distributed},
      \eqref{eq:hard} with the max-silence backstop, no freshness term.
\item \textbf{Causal-v3} --- the proposed policy,
      Algorithms~\ref{alg:sender}--\ref{alg:receiver}.
\end{itemize}

An acausal \textbf{oracle AoI} variant and three intermediate causal
variants appear in the design-evolution ablation only.

\subsection{Common random numbers}

Sharing a random seed is not the same as sharing a realisation: each policy
calls the random number generator a different number of times, so two
methods at the same seed desynchronise on the first transmission and stop
being comparable. Every stochastic input is therefore \emph{pre-drawn} as
a realisation indexed by $(\text{link},\text{timestep})$ or by physical
time, and consumed by lookup:

\begin{center}
\begin{tabular}{@{}ll@{}}
\toprule
Realisation & Indexed by \\
\midrule
forward channel loss and jitter & link, timestep \\
reverse channel loss and jitter & link, timestep \\
transmission phase & physical sender, payload class \\
link-fault pattern & graph \\
node-blackout selection & graph \\
external-force trace & physical time \\
estimator-noise trace & physical time \\
\bottomrule
\end{tabular}
\end{center}

A method that stays silent consumes nothing and disturbs nothing. The
consequence is that two methods transmitting on the same link at the same
instant meet exactly the same channel, which is what makes a
\emph{paired} comparison valid. This is verified rather than asserted: all
four methods report the same forward-realisation hash at a given seed
despite transmitting substantially different numbers of packets, and every
hash is checked against a registry computed before the sweep ran.

\subsection{Holdout validation}

The final validation used $50$ seeds drawn from a block never used during
development, with an automated assertion that no seed appears in any
seed family of the preceding experiments. The seed alone determines all
realisations, so the three network conditions at one seed share one set of
channel uniforms and differ only in the thresholds applied to them ---
common random numbers across conditions, which is what makes a
condition-to-condition difference attributable to network quality rather
than to the draw.

The final matrix contains \MetricTotalCells{} point-scenario cells over
eight selected operating points spanning $N \in \{5,20,50\}$, both plants,
ACK impairment, permanent link removal, node blackout, plant mismatch and
estimator noise, for \MetricRuns{} runs in total. Each robustness point is
a point that \emph{already exposed a limit} in its source experiment, not
the most favourable cell of that experiment.

\subsection{Statistical conventions}

These are applied uniformly and are stated here so that no reader has to
infer them from a table.

\paragraph{Key claims (holdout).} Reported as a \emph{paired} mean
difference with a $95\%$ confidence interval over $n=\MetricSeeds{}$
matched pairs,
\begin{equation}
\bar{d} \pm t_{0.975,\,n-1}\, \frac{\sigma_d}{\sqrt{n}},
\qquad d_s = m_A(s) - m_B(s) .
\label{eq:pairedci}
\end{equation}
An interval containing zero downgrades the claim to \emph{not
distinguishable at this sample size}. Adding seeds after inspecting an
interval is forbidden, because that is choosing the sample size by the
result. Where a pair is unusable --- an arm diverged, so there is no finite
error to difference --- the usable pair count is reported and the claim is
downgraded if it falls short of the pre-registered count. No pair was lost
in the results below.

\paragraph{Development experiments.} Reported as mean $\pm$ standard
deviation over $n=20$ seeds.

\paragraph{Binary safety.} Reported as \emph{unsafe / eligible} together
with the percentage, never as a percentage alone. The denominator is not
always the seed count, and hiding it would misrepresent the evidence.
Two source-specific eligibility rules are preserved:

\begin{itemize}
\item under permanent link failure, the denominator is the seeds whose
      active graph stayed connected --- a seed whose graph fell apart is a
      connectivity fact, not evidence about a method;
\item under node blackout, \emph{matched no-fault eligibility} applies: a
      seed counts only if the same method, condition and seed was already
      safe \emph{without} the fault. A configuration that was unsafe to
      begin with carries no information about blackouts.
\end{itemize}

The second rule has a consequence we report explicitly rather than let
appear as a zero: where a method was unsafe in every no-fault seed of a
cell, its eligible denominator is empty and its blackout safety is
\emph{unevaluable}, not perfect.

\paragraph{No $p$-values.} The confidence intervals answer the questions
the claims ask, and adding a significance test on top would invite reading
a threshold crossing as a result. Where an interval contains zero we say
so.

\subsection{A note on two rate normalisations}

The campaign contains two different rate conventions, and they must not be
compared:

\begin{itemize}
\item \textbf{Per channel.} Table~\ref{tab:evolution} reports the DATA
      rate on one directed channel. P10 reads exactly $10.00$~Hz.
\item \textbf{Swarm total.} Tables~\ref{tab:nominal} and
      \ref{tab:paired} report the summed rate over all directed channels.
      P10 reads \MetricDataPtenClean{}~Hz at $N=5$.
\end{itemize}

The two eras additionally used different channel counts at the same swarm
size, because the later graph convention removes in-links to the leader
(the leader reads its reference directly). Absolute traffic is therefore
comparable \emph{within} a table and not across the two conventions. Every
caption states which convention it uses.

\subsection{Reproducibility}

The campaign is frozen at a tagged commit. A single command re-runs the
test suite, re-verifies every realisation hash in a fresh process,
re-checks bit-identity between serial and parallel execution, rebuilds
every table and figure from persisted files, and writes an environment
manifest. A separate check clones the repository and reproduces the test
suite, the aggregate tables and one re-simulated seed bit-identically.
Details are in \texttt{paper/REPRODUCIBILITY.md}.
